\begin{table}[H]
\centering
\begin{tabular}{llrrrrr}
\toprule
Benchmark & Argument & $d$ & $B_{\mathrm{tr}}$ & $n_{\mathrm{tr}}$ & $n_{\mathrm{tr}}/d$ & Budget \\
\midrule
Two-state & $\Delta_2$ & 1 & 248 & 12 & $12.0$ & $1\,300$ \\
Cybersecurity & $\Delta_4$ & 3 & 184 & 20 & $6.7$ & $612$ \\
Distribution planning & $\Delta_{10}$ & 9 & 549 & 11 & $1.2$ & $2\,800$ \\
Targeted advertising & $\Delta_2$ & 1 & 248 & 12 & $12.0$ & $1\,300$ \\
Linear--quadratic & mean & 1 & 201 & 20 & $20.0$ & $4\,420$ \\
Portfolio & mean & 1 & 311 & 200 & $200.0$ & $5\,110$ \\
\bottomrule
\end{tabular}
\caption{Split of the matched simulator budget between the main trajectories $B_{\mathrm{tr}}$ and the auxiliary sensitivity trajectories $n_{\mathrm{tr}}$ of the fixed-scale transport estimator, with $d$ the number of free coordinates of the population argument. The equal-budget rule constrains $T(B_{\mathrm{tr}}+n_{\mathrm{tr}})$ but not the split, which was reallocated towards the auxiliary estimate on the continuous-state benchmarks and left at the default on the finite-state ones. The resulting auxiliary sample size per coordinate spans more than two orders of magnitude, and distribution planning is the one benchmark where it falls close to one.}
\label{tab:auxiliary-budget}
\end{table}
